\section{Ablations And Error Analysis}

\begin{table}[h]
\centering
\small
\begin{tabular}{lrrr}
\toprule
Benchmark / metric & Full router & Calibrated router & Router without delayed utility \\
\midrule
Frozen PopQA overall quality & 0.178 & 0.178 & 0.082 \\
Bundled update answer quality & 0.867 & 0.867 & 0.733 \\
Bundled update stale-answer rate & 0.133 & 0.133 & 0.267 \\
Public main answer quality (base) & 0.582 & 0.842 & 0.969 \\
Public main answer quality (small) & 0.564 & 0.844 & 0.964 \\
Public sequence answer quality (base) & 0.751 & 0.766 & 0.977 \\
Public sequence answer quality (small) & 0.740 & 0.787 & 0.974 \\
\bottomrule
\end{tabular}

\caption{Cross-benchmark ablation summary. Delayed utility helps in recurring-memory and bundled controlled-update regimes, but hurts on the public changing-answer track unless it is calibrated.}
\label{tab:arxiv-ablation}
\end{table}

Table~\ref{tab:arxiv-ablation} captures the main ablation story across regimes. On frozen PopQA, the recurring-memory result depends on \texttt{READ\_MEMORY}, \texttt{WRITE\_MEMORY}, and delayed utility. On the bundled freshness benchmark, the positive six-action result depends on \texttt{FAST\_ADAPT}, \texttt{CONSOLIDATE}, and delayed utility. On the public FreshQA-derived main slice, the direction flips: removing delayed utility dramatically improves quality, while removing temporary adaptation still hurts. That contrast is the clearest evidence that the controller does not need wholesale redesign; it needs regime-sensitive delayed-utility calibration.

\paragraph{PopQA audit.}
The PopQA audit compares a frozen pre-fix bundle against the final post-fix bundle. Before the guardrail fix, the dominant non-recurring failure was \emph{low-confidence case should have forced retrieval}. After the fix, the remaining failures were mostly \emph{retrieval was still the right choice}. This matters because it shows that the final PopQA result is not hiding a routing pathology.

\paragraph{Bundled freshness audit.}
The frozen freshness slice has only two failures, so the current manual audit is exhaustive rather than sampled. The two remaining categories are:
\begin{itemize}[leftmargin=1.5em]
\item \emph{retrieval evidence insufficient}: a volatile update where retrieval was the right action class but the visible evidence remained stale,
\item \emph{should have adapted}: a volatile update where the controller fell back to \texttt{PARAM\_ONLY} rather than using temporary adaptation.
\end{itemize}

\paragraph{Public FreshQA-derived audit signal.}
The public track now includes a hand-labeled 100-example audit comparing the full router against the calibrated router on the same \texttt{flan-t5-base} public-main slice. The dominant failure label is \emph{delayed-utility overshoot} (`74\%` of the audited failures), followed by \emph{retrieval was correctly dominant} (`26\%`). This is exactly the example-level pattern the aggregate tables suggest: the public-track failure is not primarily stale memory or premature consolidation, but over-aggressive non-retrieval behavior on a benchmark where retrieval is usually the safest action class. The paired audit matters because the calibrated router recovers many of those same cases without changing the controller action inventory.

Taken together, these audits support a consistent interpretation across the paper: the controller works in the recurring-memory and controlled-update regimes, while the public changing-answer regime primarily exposes a calibration problem that can be mitigated, but not fully erased, by regime-aware delayed-utility gating.
